\section{Introduction}
\label{sec:intro}

\emph{The picture.}
An observer looks at a periodic pattern through a lens of finite resolution $\sigma$.
The finer the pattern, the less of it survives the blur.
Paper~F~\cite{F} measured how much ``information about the direction of time'' a window keeps, through the quantity $g$.
It showed that the smallest gap $g_{\min}$ falls like $1/\sqrt{\ell}$, $\ell=\ln(1/\sigma)$, with the coefficient one, for windows with at most twenty harmonics or with decay exponent $K\ge0.003$ (which includes every window with $x\ge0.003$).
Precisely, $g_{\min}\sqrt\ell$ obeys explicit two-sided bounds for $10^{-100000}\le\sigma\le10^{-24}$, its limit lies between $1-\bar\varepsilon$ and $1$ on every class with remainder bound $\bar\varepsilon$, and it tends to one along an explicit comb.
In the language of this paper, the coefficient one is \emph{the exchange rate between fineness and blur}.
Every factor $e$ gained in fineness costs the same amount of brightness under the blur, and what remains is a fixed fraction $1/\sqrt{\ell}$.

\emph{What is new.}
Paper~F needed either a bound on the number of harmonics ($N\le20$) or a lower bound on their damping ($K\ge0.003$).
A window with many harmonics can hide its ripple in two ways.
It can spread its energy thinly over a wide band, so that the sup $\eta$ of the ripple stays small.
Or it can put a coherent block at low frequency, so that $\eta$ is large but the energy is tiny.
The second way is the reason why a conjecture we used as a working hypothesis is false (Sec.~\ref{sec:open}).
That conjecture was $\eta e^{\mu/2}\le C$, where $\mu$ measures the frequency of the ripple energy.
We built explicit windows with $\log_{10}(\eta e^{\mu/2})$ up to $32.0$ (App.~\ref{app:windows}).
The coefficient one survives on all of them; only the old tool (the $\eta$ sandwich) failed.

\emph{How we prove it.}
\begin{itemize}
\item Lemma~R and Lemma~A4* (Sec.~\ref{sec:lemmaR}) give a lower bound of $g$ for every state.
  All remainders are explicit, and the remainder $a_4$ of the curvature step (Lemma~23 of~\cite{F}) becomes relative to the ripple.
\item Lemma~S (Sec.~\ref{sec:sandwich}) replaces the sup of a low-frequency block by its energy.
\item Cell-based computer-assisted proofs (Sec.~\ref{sec:cert}) show that this lower bound exceeds the target.
\item Lemma~BE (Sec.~\ref{sec:bedge}) closes the small-$b$ edge by measuring the curvature cross term in mean square instead of in sup.
\item The boxes and Lemma~BE cover every state of the theorem; this is shown line by line and checked on $10^5$ random states, a check that also finds every deliberately removed box (Sec.~\ref{sec:cert}).
\end{itemize}
Table~\ref{tab:status} lists every result with its status.
Throughout we keep the distinction of~\cite{F}: proof, computer-assisted proof, derivation, numerical, conjecture.

\begin{table}
\caption{Results of this paper and their status.}
\label{tab:status}
\begin{ruledtabular}
\small
\begin{tabular}{p{0.50\columnwidth}p{0.40\columnwidth}}
result & status\\ \hline
Lemma R (lower bound for every state) & proof (Sec.~\ref{sec:lemmaR})\\
Lemma A4* ($a_4$ relative) & proof\\
Lemma H (high part and $\theta_\tau$ kernels) & proof (App.~\ref{app:tau})\\
$\tau$ cells, factor $0.99977$ on $[20,80]$ & computer-assisted (App.~\ref{app:tau})\\
$\tau$ factors decrease on $\tau\ge20$ & proof (App.~\ref{app:tau})\\
Lemma S, Lemma S$'$, Corollary S & proof (Sec.~\ref{sec:sandwich})\\
$x$-free envelope (Lemma Z3$'$) & proof (Sec.~\ref{sec:envelope})\\
main reduced inequality, $\ell\le6907.8$, with the $\mu$ tails and $b\to0$ & computer-assisted (Sec.~\ref{sec:cert})\\
full Lemma R* bound, $\tau\ge20$, $b\ge10^4\sigma^2$, $\ell\le960.8$, with the $\mu$ tails and the wedge & computer-assisted (Sec.~\ref{sec:cert})\\
covering of the states by the boxes and Lemma BE & proof (line by line); random check of $10^5$ states, numerical (Sec.~\ref{sec:cert})\\
Lemma BE (every state with $b<10^4\sigma^2$) & proof (Sec.~\ref{sec:bedge})\\
seventeen thin-block windows & computer-assisted (Sec.~\ref{sec:cert})\\
Theorem H1 (every state with $\tau\ge20$) & computer-assisted, modulo (P6) (Sec.~\ref{sec:theorem})\\
envelope values ($1.7116$ at $10^{-24}$, $1.0063$ at $10^{-3000}$) & numerical\\
conjecture C ($\eta e^{\mu/2}\le C$) & false (App.~\ref{app:windows})\\
local corrections (L1)--(L3), Lemma A* & proof (Sec.~\ref{sec:open:tau})\\
local Z3$'$ with one universal constant & false; ratio $26.9$ and $71.4$ on explicit packets (Sec.~\ref{sec:open:tau})\\
Lemma CM-1 (local brightness) & proof (Sec.~\ref{sec:open:tau})\\
Lemma CN-1 and the joint bounds $\Psi_2$, $\Psi_2^*$ & proof (Sec.~\ref{sec:open:tau})\\
Lemma CN-2 (local $F$ term) & proof (Sec.~\ref{sec:open:tau})\\
Lemma R$_\theta$ (local route, explicit remainders) & proof (Sec.~\ref{sec:open:tau})\\
sharp joint bound $\Psi\min(1,\sqrt{2\kappa})$ & conjecture; survives the search (App.~\ref{app:joint})\\
local envelope with $\Psi_2^*$ above $T_0$ ($\ell\le960.8$) & numerical (Sec.~\ref{sec:open:tau})\\
states with $\tau<20$ & open (Sec.~\ref{sec:open:tau})\\
full bound for $\ell>960.8$, limit $\sigma\to0$ & open (Sec.~\ref{sec:open:sigma})\\
\end{tabular}
\end{ruledtabular}
\end{table}

\emph{Picture.} The lens sees a fine ripple only through its blurred brightness.
This paper shows that the brightness that survives does not depend on how many harmonics make up the ripple, only on how much energy they carry at each frequency.
